\documentclass[11pt,a4paper]{article}
\usepackage[T1]{fontenc}
\usepackage{lmodern,amsmath,amssymb,amsthm,microtype,booktabs}
\usepackage[margin=25mm]{geometry}
\usepackage[colorlinks=true,urlcolor=blue,linkcolor=blue]{hyperref}
\newcommand{\dd}{\mathop{}\!\mathrm d}
\newcommand{\sech}{\operatorname{sech}}
\newcommand{\HH}{\mathcal H}
\title{Group pairing and a physical two-point form\\in a local dS$_3$ observer model}
\author{Alan Day \and Chengheng\thanks{AI research assistant (OpenAI).}}
\date{Working research note --- 31 October 2026}
\begin{document}\maketitle
\begin{abstract}
We fix the group-averaged norms of the vacuum and one-matter sectors
in the previously specified local, zero-gravity dS$_3$ model.
A linear left--right gluing condition selects a normalizable mass wavefunction.
The invariant one-matter norm then gives a positive, finite-mass two-point
spectral measure, including the full angular invariant sector.
These are normalization-specific derivations using known group theory.
They do not construct the full Fock algebra or establish finite-gravity holography.
\end{abstract}
\section{Model and the two-observer pairing}
Set the radius $\ell=1$ and Newton's constant $G_N=0$. Use the proper group
$\mathsf G=SO^+(3,1)$ and two distinguishable scalar BD observers
with mass labels $\nu\ge s_*>0$. Matter is one fixed free principal-series
scalar with label $p>0$, in its BD vacuum. The local insertion has coupling one.
The previous spacetime normalization is
\[
v_{\nu lm}=\sqrt{\mu(\nu)}u_{\nu lm},\qquad
\mu(\nu)=\frac{\nu\tanh\pi\nu}{\pi};
\]
matter modes remain KG-normalized. This $\mu$ is not a group density.

Give $K=SO(3)$ Haar probability measure and fix the overall Haar constant by
\begin{equation}
 \int_{\mathsf G}F(g)\dd g=
 \int_0^\infty\sinh^2r\,\dd r\int_K\dd k_1\dd k_2\,F(k_1a_rk_2).
\end{equation}
The normalized spherical function and its positive-label orthogonality are
\begin{align}
 \phi_\nu(r)&=\frac12\int_{-1}^1(\cosh r-x\sinh r)^{-1-i\nu}\dd x
 =\frac{\sin\nu r}{\nu\sinh r},\\
 \int_{\mathsf G}\phi_a\phi_b\,\dd g
 &=\frac{\pi}{2a^2}\delta(a-b),\qquad
 K_2(\nu)=\frac{\pi}{2\nu^2},\quad \rho_{\mathsf G}(\nu)=\frac{2\nu^2}{\pi}.
\end{align}
The sine-transform identity fixes the constant. The same dS$_3$ density
already appears in Yang--Zhang--Zheng \cite{YYZ}, (2.23)--(2.25) and Appendix A.
Our calculation is a normalization check, not a new Plancherel formula.
Positive labels choose one reflection representative.

\newpage
\section{Rigging coefficients and the three-body norm}
In the BD $K$ basis the invariant two-body distribution is
\begin{equation}
 I_\nu=\sum_{l,m}(-1)^{l+m}|\nu,lm\rangle_L|\nu,l,-m\rangle_R .
\end{equation}
Indeed $\Theta|lm\rangle=(-1)^{l+m}|l,-m\rangle$, with conjugation
of coefficients, is an antiunitary intertwiner with $\Theta^2=1$.
Rotations obey the usual singlet identity. For boosts, the real adjacent-$l$
matrix elements and the $l$ parity give the required sign under conjugation.

Start with seeds smooth and compact in both mass labels, away from zero
and the cutoff, with finitely many $K$ types. Perform mass integration
before group integration. Compact-mass Harish--Chandra expansions and
integration by parts suppress the leading $e^{-2r}$ two-body term by
arbitrary inverse powers of $r$; remaining terms decay faster.
Thus the smeared pairing is absolutely integrable.
Fixed-label pairings are distributions, not ordinary integrals.
Finite $K$ type is a convenient dense test domain; its group-translate span
may be used for an invariant domain \cite{MM,MY}.

If $c_\psi(\nu)=I_\nu^\dagger\psi(\nu,\nu)$, group averaging gives
$f=K_2c_\psi$ and
\begin{equation}
 \langle\eta\psi,\eta\chi\rangle
 =\int_{s_*}^\infty K_2\overline{c_\psi}c_\chi\,\dd\nu
 =\int_{s_*}^\infty\rho_{\mathsf G}\bar f g\,\dd\nu .
\end{equation}
Accordingly the coefficient of the formal invariant distribution
$\Psi_f=\int f(\nu)I_\nu\dd\nu$ has norm $\int\rho_{\mathsf G}|f|^2$.
This inverse relation is essential: neither a squared delta function
nor an interchange of $K_2$ and $\rho_{\mathsf G}$ is allowed.

For three fixed positive labels $(a,b,p)$, finite-$K$ matrix coefficients
have magnitude $O((1+r)^N e^{-r})$ per particle. The product is integrable
against $\sinh^2r\dd r$. The scalar trilinear invariant is unique in this
generic principal-series domain \cite{DS}. Normalize its distribution
$I_{abp}$ to have $000$ component one. Write
\[
 S=a+b+p,\qquad D=\cosh\frac{\pi S}{2}
       \prod_{n=a,b,p}\cosh\frac{\pi(S-2n)}{2}.
\]
Then the rank-one rigging form is $K_3|I_{abp}\rangle\langle I_{abp}|$, where
\begin{align}
 K_3&=\int_0^\infty\sinh^2r\,\phi_a\phi_b\phi_p\,\dd r\\
 &=\frac{\pi}{8abp}\left(\sum_{n=a,b,p}\tanh\frac{\pi(S-2n)}2
                   -\tanh\frac{\pi S}2\right)
 =\frac{\pi}{8D}\prod_{n=a,b,p}\frac{\sinh\pi n}{n}>0 .
\end{align}
Use $\int_0^\infty\sin(kr)/\sinh r\,\dd r=(\pi/2)\tanh(\pi k/2)$.
An invariant coefficient $h(a,b)I_{abp}$ therefore has physical norm
$\int |h|^2/K_3\,\dd a\dd b$. The $000$ component fixes normalization;
it does not truncate the matter harmonics.

\newpage
\section{Linear gluing and a normalizable state}
The previous local vertex, with outgoing observer $a$, incoming observer
$b$ and emitted matter $p$, has spherical entry
\begin{align}
 B_*&=\frac{\pi}{4\sqrt{4\pi}\,D}
          \sqrt{\prod_{n=a,b,p}\frac{\sinh\pi n}{n}},\\
 E_{abp,000}&=\sqrt{\mu(a)\mu(b)}B_*\cosh\pi b .
\end{align}
Left emission from $\Psi_f$ has invariant coefficient $f(b)E_{abp,000}$;
right emission has $f(a)E_{bap,000}$. Both multiply the same $I_{abp}$.
At other angular labels the singlet phases must be retained.

Equality of these two linear scalar insertions is equivalent to
\begin{equation}
 f(b)\cosh\pi b=f(a)\cosh\pi a,\qquad
 f_{\rm lin}(\nu)=\sech\pi\nu
\end{equation}
up to an overall constant (almost everywhere on the connected mass domain).
Kolchmeyer--Liu already obtained this function in dS$_2$ \cite{KL}.
Here it follows again from the declared dS$_3$ vertex and norms.
We call it the \emph{linear-gluing state}, or a candidate HH state.
Gluing for arbitrary composites and the complete observable algebra has
not been established by this one-particle calculation.

The normalization and probability density are
\begin{equation}
 Z(s)=\int_s^\infty\frac{2\nu^2}{\pi}\sech^2\pi\nu\,\dd\nu,\qquad
 \dd P(\nu)=Z(s_*)^{-1}\rho_{\mathsf G}(\nu)\sech^2\pi\nu\,\dd\nu .
\end{equation}
With $q=e^{-2\pi s}$, integration by parts gives
\begin{equation}
 Z(s)=\frac{4s^2}{\pi^2}\frac{q}{1+q}
      +\frac{4s}{\pi^3}\log(1+q)
      -\frac{2}{\pi^4}\operatorname{Li}_2(-q).
\end{equation}
The formal endpoint check is $Z(0)=1/(6\pi^2)$; the actual model retains
$s_*>0$. At large label, the unnormalized weight is
$(8/\pi)\nu^2e^{-2\pi\nu}(1+O(e^{-2\pi\nu}))$.
The mass is $\sqrt{1+\nu^2}$, so finite-label Boltzmann interpretations
must not replace mass by $\nu$ without specifying the approximation.

The full one-matter norm now removes the three-body normalization:
\begin{equation}\boxed{
 Q(a,b;p)=\frac{|f_{\rm lin}(b)E_{abp,000}|^2}{K_3(a,b,p)}
         =\frac{\mu(a)\mu(b)}{8D}.}
\end{equation}
This kernel is positive and symmetric in the two observer labels.
Unlike a finite-mode probe sum, it represents the full angular invariant
one-matter sector. Uniqueness of the invariant and independent
non-spherical group integrals support the normalization.

\newpage
\section{A closed creation block and the two-point form}
Let $\HH_0=L^2([s_*,\infty),\rho_{\mathsf G}\dd b)$ and
$\HH_1=L^2([s_*,\infty)^2,K_3^{-1}\dd a\dd b)$ in the invariant coefficient
representations. The zero-to-one creation block is
\[
 (Af)(a,b)=f(b)E_{abp,000}.
\]
Its norm is exactly
\begin{align}
 \|Af\|_1^2&=\int_{s_*}^\infty\rho_{\mathsf G}(b)|f(b)|^2
                           {\cal S}_p(b)\dd b,\\
 {\cal S}_p(b)&=\frac{\cosh^2\pi b}{\rho_{\mathsf G}(b)}
                     \int_{s_*}^\infty Q(a,b;p)\dd a .
\end{align}
Since $D\ge e^{2\pi\max(a,b)}/16$ and
$\mu(a)\mu(b)\le ab/\pi^2$,
\begin{equation}
 Q(a,b;p)\le\frac{2ab}{\pi^2}e^{-2\pi\max(a,b)},\qquad
 {\cal S}_p(b)\le\frac{b}{2\pi}+\frac{1}{2\pi^2}
                         +\frac{1}{4\pi^3s_*}.
\end{equation}
The second inequality follows by splitting the $a$ integral at $b$ and
using $\cosh^2\pi b\le e^{2\pi b}$. A finite first mass moment is sufficient.
The exact maximal domain is the displayed weighted integrability domain.
Fiberwise multiplication gives a closed densely defined map: along an
almost-everywhere convergent subsequence of a graph-convergent sequence,
the limiting fiber is still $f(b)E_{abp,000}$. Compact mass supports are
dense and have finite rate. The adjoint defines the reverse block.

This constructs two sectors and their connecting block. It is not a proof
of all particle sectors, a common Fock core, or self-adjointness of a full field.

Take $h(\nu)=\sqrt{1+\nu^2}$ as the rest-mass clock Hamiltonian of this
specified model. Evolving the output and input observer masses gives
the normalized vacuum two-point quadratic form
\begin{equation}\boxed{
 C_p(t)=\frac1{Z(s_*)}\int_{s_*}^\infty\!\dd a\dd b\,
 Q(a,b;p)e^{-it[h(a)-h(b)]}.}
\end{equation}
It is finite for every real $t$. Symmetry makes it real and even;
the positive pushforward measure under $\Omega=h(a)-h(b)$ makes it
positive-definite, which does not imply pointwise positivity.
At radius $\ell$, the physical frequency is $\Omega/\ell$ and, with the
corresponding canonical scalar normalization, $C$ has units $\ell^{-1}$.

For a square cutoff $L>s_*$, the absolute unnormalized outer-tail bound is
\begin{equation}
 \frac{2e^{-\beta L}}{\pi^2}
 \left[\frac{L^3-s_*^2L}{\beta}
       +\frac{3L^2-s_*^2}{\beta^2}+\frac{6L}{\beta^3}+\frac6{\beta^4}\right],
 \qquad \beta=2\pi .
\end{equation}
It is uniform in real time. Divide by $Z$ for the normalized form.
No small tail alone certifies the interior oscillatory quadrature.

\newpage
\section{Comparison with the earlier recoil ansatz}
This is a comparison of a derived local $G_N=0$ model with the previous
scalar recoil ansatz. It is not a finite-$\kappa$ insertion dictionary.
Set
\[
 M=(a+b)/2,\quad\omega=a-b,\quad
 A_0=\cosh2\pi M,\quad B_0=\cosh\pi\omega,\quad C_0=\cosh\pi p.
\]
The product identity gives $D=(A_0+C_0)(B_0+C_0)/4$.
Recall the earlier definitions
\begin{align}
 k_p(\omega)&=\frac1{8\pi(B_0+C_0)},\\
 R_{\rm old}&=\left(1-\frac{\omega^2}{4M^2}\right)
                   \frac{A_0-B_0}{A_0+C_0},\\
 Q_{\rm old}(a,b;p)&=M^2e^{-2\pi M}k_p(\omega)R_{\rm old}.
\end{align}
Direct algebra in the present norms yields
\begin{align}
 \frac{Q(a,b;p)}{\rho_{\mathsf G}(M)\sech^2\pi M}
       &=k_p(\omega)R_{\rm phys}(M,\omega;p),\\
 R_{\rm phys}&=R_{\rm old}\frac{A_0+1}{A_0+B_0},\qquad
                  0<R_{\rm phys}\le1 .
\end{align}
For fixed $(\omega,p)$, $R_{\rm phys}\to1$ as $M\to\infty$.
With the common overall scale calibrated at large mass, an equivalent
comparison is
\begin{equation}\boxed{
 \frac{\pi}{8}Q=Q_{\rm old}\,n_{\rm loc}(a,b),\qquad
 n_{\rm loc}(a,b)=
 \frac1{(1+e^{-2\pi a})(1+e^{-2\pi b})}.}
\end{equation}
Previously, leading matching allowed a family of finite-mass completions.
The local insertion, group norms and chosen state now select this completion
\emph{in this model}. They do not remove the missing state, inner-product
and contour inputs of a finite-$\kappa$ proposal. In particular
$n_{\rm loc}(M,M)\ne1$ at finite $M$: an earlier diagonal convention
was a diagnostic choice, not a physical constraint.

The label frequency is not the proper clock energy:
\begin{equation}
 |h(a)-h(b)-(a-b)|\le\frac{|\omega|}{2\min(a,b)^2}.
\end{equation}
To see this, integrate $1-h'(x)=1-x/\sqrt{1+x^2}\le1/(2x^2)$
between the two masses. For bounded times, the phase error is at most
$|t|$ times this bound. Fixed-label spectral convergence alone does not
give all-time convergence of the proper-time Fourier transform.
No such all-time relative estimate is claimed here.

\newpage
\section{Checks, reproducibility and next question}
The saved main diagnostics contain 81 checks; a separate run contains
16 checks of non-spherical group averaging. Main arithmetic uses
50 decimal digits and the independent run uses 45; time quadratures use
NumPy double precision with Gauss--Legendre orders 80 and 120.
The code, inputs, versions and raw complex values are retained.

\begin{center}
\begin{tabular}{lr}
\toprule
Main diagnostic group & Checks\\ \midrule
Spherical angular representation &9\\
Two-body compact mass packet and analytic tail &2\\
Three-body direct radial integral, product and tail &9\\
Higher-angular left--right gluing &15\\
Vertex/group quotient, symmetry and recoil ratio &9\\
State normalization and creation-rate bound &6\\
Mean-label spectrum and physical energy bound &15\\
Time quadrature refinement and imaginary part &16\\ \bottomrule
\end{tabular}
\end{center}
The triple samples are $(.7,1.1,.6)$, $(1.4,.9,1.2)$ and $(3.2,.4,.5)$.
The direct triple integral ends at 65 with tail
$2\operatorname{atanh}(e^{-65})/(abp)$.
The pair test uses $g(\nu)=(\nu-.4)^4(1.6-\nu)^4$ on $[.4,1.6]$,
mass coefficient $\nu g(\nu)$, and radial cutoffs 40 and 80.
Its Schur target is $(\pi/2)(1.2)^{17}B(9,9)$.
Parseval applied to $g^{(4)}$ gives an $R^{-8}$ tail bound.

The independent angular test uses the rotational singlet
\[
 S_l=(2l+1)^{-1/2}\sum_m(-1)^{l-m}|lm\rangle|l,-m\rangle|00\rangle .
\]
For $l=1,2$ and the first two triples it compares the actual group
integral to $K_3 I_{abp}(S_l)$, where
$I_{abp}(S_l)=\sqrt{2l+1}\,T_{ll0}/T_{000}$ in the previous vertex convention.
Boosts are computed from compact-boundary angular integrals and
independent finite polynomial primitives. The $r>70$ tail is bounded by
\[
 \frac{2\sqrt{2l+1}}{p(1-e^{-140})}\,e^{-70}(70^2+140+2).
\]
This tests non-spherical components of the rank-one form, not only its
normalization component.

\begin{center}
\begin{tabular}{ccc}
\toprule
$t$ & $C_{.6}(t)$, $s_*=.2$ & $C_{1.2}(t)$, $s_*=.8$\\ \midrule
0   & .005995878862 & .001281989885\\
.5  & .005942866399 & .001261299720\\
2   & .005231147771 & .001002967854\\
5   & .002966801879 & .000418357992\\ \bottomrule
\end{tabular}
\end{center}
The respective normalized outer-tail bounds at $L=s_*+8$ are
$5.11\,10^{-20}$ and $8.94\,10^{-21}$; interior refinement errors are larger.
These are empirical quadratures with analytic outer tails, not interval
certificates of total error.

\newpage
\section{Scope and continuation}
The next calculation is the two-matter sector and its physical four-point
recoupling, including all principal-series spin channels. It must establish
the vacuum-to-one-to-two domain, keep each observer projection and
recover the independent-free-probe disconnected benchmark in its actual
large-mass regime. For two independent probe labels $p,p'$ that target is
\[
 \frac{pp'}{16\pi^2\sinh\pi p\,\sinh\pi p'}
\]
in the previously declared antipodal convention; equal species have
additional Wick terms. This benchmark has not been checked by the
present two-point run.

Three distinct open questions remain: composite gluing and a full Fock
observable algebra; the finite-$\kappa$ inner product and sewing contour;
and finite-$G_N$ dynamics and a complete holographic dual.
Our derived $n_{\rm loc}$ answers none of these on its own.

\paragraph{Local project record.}
The Chinese derivation is \texttt{notes/stage4-ds3-group-pairing.md}.
Raw diagnostics are in \texttt{results/stage4/ds3-group-pairing/}.
The model implementation is \texttt{code/ds3\_group\_pairing.py};
the two independent validation drivers and the build script are in
\texttt{code/}. This is a private working note, not a published
originality claim.

\begin{thebibliography}{9}
\bibitem{YYZ}
Z. Yang, Y. Zhang and W. Zheng,
\emph{Comments on the de Sitter Double Cone},
\href{https://arxiv.org/html/2505.08647v1}{arXiv:2505.08647v1},
Sec.\ 2.2 and Appendix A. The cached v1 title is used here.
\bibitem{KL}
D. Kolchmeyer and H. Liu,
\emph{Quantum observers in de Sitter space},
\href{https://arxiv.org/html/2411.08090v1}{arXiv:2411.08090v1},
Secs.\ 4 and 6. The linear-gluing construction is a direct dS$_2$ predecessor.
\bibitem{MM}
D. Marolf and I. A. Morrison,
\emph{Group Averaging for de Sitter free fields},
\href{https://arxiv.org/html/0810.5163v1}{arXiv:0810.5163v1},
Sec.\ 2.3 and boost matrix elements.
\bibitem{MY}
D. Marolf and Z. Yu,
\href{https://arxiv.org/html/2610.24073v1}{arXiv:2610.24073v1},
Secs.\ 5.1--5.2, compact-mass group averaging and positivity.
\bibitem{DS}
S. E. Derkachov and V. P. Spiridonov,
\emph{The 6j-symbols for the SL(2,C) group},
\href{https://arxiv.org/html/1711.07073v2}{arXiv:1711.07073v2},
scalar intertwiners and reflection conventions.
Their raw CG measure is not imported into our Haar normalization.
\end{thebibliography}
\end{document}
